% Part: incompleteness
% Chapter: arithmetization-syntax
% Section: introduction

\documentclass[../../../include/open-logic-section]{subfiles}

\begin{document}

\olfileid{inc}{art}{int}
\olsection{مقدمة}

وصل المنطق بقابلية الحساب يقتضي أن نصوغ الحديث
عن رموزه وحدوده وال!!{formula}s وال!!{derivation}s،
وعن خواصها وعلاقاتها والعمليات عليها، صياغة تصلح
للحساب. ويمكن في بعض النماذج تناول هذه الأشياء
مباشرةً بدوال وعلاقات قابلة للحساب على الرموز
ومتتالياتها وما يبنى منها؛ فكائنات التركيب
المنطقي منتهية، ورموزها مأخوذة من مجموعات
!!a{enumerable}. غير أن نماذج أخرى، كالدوال
العودية، لا تتناول إلا الأعداد ودوالها وعلاقاتها.
ثم إن غايتنا معالجة التركيب داخل نظريات معينة،
ولا سيما نظريات لغة الحساب. فهنا لا بد من
\emph{حوسبة} التركيب: تمثل الكائنات التركيبية
بأعداد، والعمليات بدوال حسابية، والعلاقات
بعلاقات حسابية. وأصل هذه الفكرة، أي إسناد
الأعداد إلى الكائنات التركيبية، يرجع إلى لايبنتس.

وترميز الرموز المفردة بالأعداد أمر قريب.
وفيه شيء من الصعوبة؛ فال!!{variable}s غير
متناهية، وكذلك ال!!{function}s لكل رتبة~$n$.
إلا أن تعيين الرموز العددية على نظام ممكن،
بحيث يكون، مثلًا، لـ$\Obj v_2$ و$\Obj v_3$
رمزان متباينان. وأما المتتاليات، ومنها الحدود
وال!!{formula}s، فأشد صعوبة. فإذا أمكن
ترميز متتاليات الأعداد ومعالجتها حسابيًا،
كترميزها بقوى الأعداد الأولية، امتد ترميز
الرموز المفردة إلى متتالياتها، ثم إلى
متتاليات ال!!{formula}s وسائر ترتيباتها،
كال!!{derivation}s. وهذا الترميز الموسع
يسمى «ترقيم غودل»، ويعيّن لكل حد
و!!{formula} و!!{derivation} عدد غودل.

فإذا رمزنا متتالية الرموز بمتتالية رموزها العددية،
واخترنا للمتتاليات ترميزًا تقبل معالجته الدوال
القابلة للحساب، أمكن معالجة أعداد غودل
بتلك الدوال أيضًا. بل ستكون الدوال التي
نحتاج إليها عودية بدائية. فاستخراج طول
المتتالية وعنصرها رقم $i$ من رمزها العددي
عمليتان عوديتان بدائيتان. وإذا كان الرمز
عدد غودل ل!!a{formula}~$!A$، تبين فورًا
إمكان حساب طول !!a{formula} ورمز علامتها
رقم $i$ في !!a{formula} من عدد غودل لـ~$!A$.
وأشق قليلًا إثبات أن كون العدد رمزًا لحد
حسن التكوين أو ل!!{derivation} صحيح خاصية
عودية بدائية. إلا أن هذا ممكن، لأن
المتتاليات التي نعنيها، أي الحدود
وال!!{formula}s وال!!{derivation}s،
معرّفة تعريفًا استقرائيًا.

ولنوضح ذلك بالاستبدال. لتكن $!A$ صيغة،
و$x$ متغيرًا، و$t$ حدًا. فـ$\Subst{!A}{t}{x}$
حاصل إحلال~$t$ محل كل ورود حر لـ~$x$
في~$!A$. ولنفرض أن أعداد غودل لـ$!A$
و$x$ و$t$ هي $k$ و$l$ و$m$ بالترتيب.
وبالترقيم نفسه يكون لـ$\Subst{!A}{t}{x}$
عدد غودل، ولنسمه~$n$. فالدالة التي تأخذ
$k$ و$l$ و$m$ إلى $n$ هي نظير الاستبدال
في الحساب. فكما يأخذ الاستبدال $!A$
و$x$ و$t$ إلى $\Subst{!A}{t}{x}$،
تأخذ دالته المحوسبة أعداد غودل $k$
و$l$ و$m$ إلى~$n$. وسنثبت أنها عودية بدائية.

ولا تنحصر فائدة حوسبة التركيب في هذا الاعتبار
المجرد، على أهمية اكتشاف أن لغة الحساب،
مع خلوها من أدوات خاصة للكلام على اللغات،
تستطيع التعبير عن خواص معقدة للتعبيرات.
فمنه ننتقل إلى السؤال عما \emph{تثبته}
نظرية، كحساب بيانو، عن لغتها هي، ومن
ذلك قابلية ال!!{sentence}s للبرهنة أو
صدقها. ومن هنا تأتي مبرهنات غودل في
حدود البرهنة وتارسكي في امتناع تعريف الصدق.
وللحوسبة أثر كذلك في نظرية قابلية الحساب:
فبها تثبت نتائج عن القوة الحسابية للنظريات
وعلاقات نماذج الحساب بعضها ببعض.
وهي أيضًا مثال لحوسبة كائنات وخواص أخرى؛
فيمكن ترميز تهيئات آلات تورنغ وحساباتها
على نحو مماثل، ثم محاكاة الحساب في نموذج،
كآلات تورنغ، بنموذج آخر، كالدوال العودية.

\end{document}
